\section{Introduction}
\label{sec:intro}

Every computation in a transformer flows through the residual stream---a sequence of hidden state vectors that carry all information between layers. A natural question follows: {\bf how much information can a single hidden state vector actually store?} The answer depends critically on how that information is encoded and decoded, yet prior work has studied only narrow slices of this spectrum. \citet{kuratov2025cramming} showed that a single 4096-dimensional vector can encode 1,568 tokens when decoded through a full frozen LLM, while \citet{elhage2022superposition} demonstrated that sparse features can be stored in superposition. Missing from the literature is a systematic, apples-to-apples comparison across the full range of decoder complexities.

This gap matters for several active research areas. Soft prompt methods \citep{belrose2023tuned,pal2023future} and memory token approaches implicitly assume that hidden states can carry substantial information, but the capacity limits under different decoding strategies remain unclear. Model compression techniques need to know how much information hidden states carry versus how much the decoder contributes. Mechanistic interpretability efforts \citep{geva2021transformer,dar2023analyzing} benefit from understanding what fraction of a hidden state's capacity is actually utilized.

We address this gap with the first systematic measurement of hidden state information capacity across five encoding schemes of increasing decoder complexity, all evaluated on the same \gptwo model ($\dmodel = 768$). Our schemes range from raw ASCII bit-packing (no decoder, 24,576 bits at fp32) through vocabulary encoding, geometric vector walks, and unembedding-based projection, to full frozen-model decoding (578 bits from a single vector, 7,762 bits from 32 vectors). Using random tokens sampled uniformly from the vocabulary---which have maximum entropy of $\sim$15.62 bits per token and prevent the model's language prior from contributing---we isolate the pure storage capacity of the hidden state under each decoding scheme.

Our experiments reveal several findings. First, raw bit-packing without any decoder stores the most information: 24,576 bits in a 768-dimensional fp32 vector. Second, the unembedding matrix can decode random tokens from as few as 3 dimensions per token with a learned linear projection (200 tokens at 100\% accuracy), but accuracy collapses at 2 dimensions per token. Third, superposition encoding---which succeeds for sparse binary features \citep{elhage2022superposition}---fails completely for dense token IDs. Fourth, a single frozen transformer layer provides almost no capacity benefit beyond position embeddings, while the full 12-layer model dramatically improves decoding. Finally, scaling from 1 to 32 \memvec vectors yields 99.4\% accuracy on 500 random tokens, with peak efficiency of 0.46 bits per learnable parameter at 16 vectors.

Our contributions are:
\begin{itemize}[leftmargin=*,itemsep=0pt,topsep=0pt]
    \item We conduct the first systematic comparison of hidden state information capacity across five encoding schemes spanning the full decoder-complexity spectrum, from zero-decoder bit-packing to full frozen-model decoding.
    \item We identify a sharp capacity threshold at 3 dimensions per token for unembedding-based decoding and demonstrate that superposition encoding fails for dense token IDs, contrasting with its success for sparse features.
    \item We quantify the decoder--capacity tradeoff: simpler encodings store more raw bits (up to 32 $\bpd$), while model-based decoders achieve only 0.75 $\bpd$ from a single vector but leverage learned structure to recover meaningful token sequences.
    \item We show that multi-vector scaling achieves 99.4\% accuracy on 500 random tokens with 32 vectors, reaching peak efficiency of 0.46 $\bpp$ at 16 vectors.
\end{itemize}

The rest of this paper is organized as follows. \Secref{sec:related} reviews related work on hidden state capacity and decoding. \Secref{sec:method} describes our five encoding schemes and experimental setup. \Secref{sec:results} presents results across all schemes, and \secref{sec:discussion} discusses the decoder--capacity tradeoff, limitations, and implications. \Secref{sec:conclusion} concludes.
